% GENERATED FILE. Edit canonical lessons, then run npm run generate:books.
% canonical-source-hash: fnv1a64:b169bc460a8ee047
% canonical-lessons: IT-C367-volontariato, IT-C367-volontario, IT-C367-associazione, IT-C367-beneficenza, IT-C367-comunita

\chapter{Dare una mano}
\label{ch:it-dare-una-mano}
\hlchaptermodality{367}

\begin{chapteropening}
\textbf{By the end of this chapter:} \emph{I can talk about helping others: volunteering, a volunteer, an association, charity and the community.}

Complete the last lesson of chapter 367: I can talk about helping others: volunteering, a volunteer, an association, charity and the community.
\end{chapteropening}

\section[il volontariato]{il volontariato — the voluntary work, volunteering}

\label{lesson:IT-C367-volontariato}



\begin{quote}
Before the new one: say the Italian for to discriminate, then the Italian for the immigration.
\end{quote}

\subsection*{You'll want to know: il volontariato}
\textbf{il volontariato} — \textquotedblleft{}the voluntary work, volunteering\textquotedblright{}.

\begin{quote}
Il sabato faccio volontariato in ospedale.
\end{quote}

\begin{quote}
\emph{On Saturdays I do voluntary work at the hospital.}
\end{quote}

\emph{Fare volontariato} — to volunteer.

\begin{cousinweb}
\emph{volontariato} is built on \emph{volontario}, a learned borrowing from Latin \emph{voluntārius} \textquotedblleft{}willing, of free will\textquotedblright{} (Lewis \& Short), from \emph{voluntās} \textquotedblleft{}will\textquotedblright{}, built on \emph{velle} \textquotedblleft{}to want\textquotedblright{}.
\end{cousinweb}

\begin{grammarlens}[title={What you've built}]
One more word for telling, explaining and describing at B1.
\end{grammarlens}

\subsection*{Your turn}
\begin{itemize}
\raggedright
  \item \emph{Say it:} \emph{il volontariato}
  \item \emph{Say it:} \emph{il volontariato}, once more
  \item \emph{Say it:} a sentence of your own with \emph{il volontariato}
  \item \emph{Recall:} say the Italian for to discriminate, then the Italian for the immigration, then say \emph{il volontariato} again
\end{itemize}

\begin{tcolorbox}[breakable,colback=teal!4,colframe=teal!35!black,title={Before you move on}]
What does il volontariato mean? (\textquotedblleft{}The voluntary work, volunteering\textquotedblright{}.) Say it once more.
\end{tcolorbox}
\section[il volontario]{il volontario — the volunteer}

\label{lesson:IT-C367-volontario}



\begin{quote}
Before the new one: say the Italian for the immigration, then the Italian for the voluntary work.
\end{quote}

\subsection*{You'll want to know: il volontario}
\textbf{il volontario} — \textquotedblleft{}the volunteer\textquotedblright{}.

\begin{quote}
I volontari portano da mangiare agli anziani.
\end{quote}

\begin{quote}
\emph{The volunteers bring food to the elderly.}
\end{quote}

A woman is \emph{la volontaria}. As an adjective: voluntary.

\begin{cousinweb}
\emph{volontario} is a learned borrowing from Latin \emph{voluntārius} \textquotedblleft{}willing, of free will\textquotedblright{} (Lewis \& Short), from \emph{voluntās} \textquotedblleft{}will\textquotedblright{}, built on \emph{velle} \textquotedblleft{}to want\textquotedblright{}.
\end{cousinweb}

\begin{grammarlens}[title={What you've built}]
One more word for telling, explaining and describing at B1.
\end{grammarlens}

\subsection*{Your turn}
\begin{itemize}
\raggedright
  \item \emph{Say it:} \emph{il volontario}
  \item \emph{Say it:} \emph{il volontario}, once more
  \item \emph{Say it:} a sentence of your own with \emph{il volontario}
  \item \emph{Recall:} say the Italian for the immigration, then the Italian for the voluntary work, then say \emph{il volontario} again
\end{itemize}

\begin{tcolorbox}[breakable,colback=teal!4,colframe=teal!35!black,title={Before you move on}]
What does il volontario mean? (\textquotedblleft{}The volunteer\textquotedblright{}.) Say it once more.
\end{tcolorbox}
\section[l'associazione]{l'associazione — the association, the organisation}

\label{lesson:IT-C367-associazione}



\begin{quote}
Before the new one: say the Italian for the voluntary work, then the Italian for the volunteer.
\end{quote}

\subsection*{You'll want to know: l'associazione}
\textbf{l'associazione} — \textquotedblleft{}the association, the organisation\textquotedblright{}.

\begin{quote}
Lavoro per un'associazione che aiuta i bambini.
\end{quote}

\begin{quote}
\emph{I work for an organisation that helps children.}
\end{quote}

\emph{Un'associazione culturale} — a cultural club.

\begin{cousinweb}
\emph{associazione} is built on \emph{associare} \textquotedblleft{}to join together\textquotedblright{}, a learned borrowing from Late Latin \emph{associāre} \textquotedblleft{}to join to\textquotedblright{} (Lewis \& Short), built on \emph{socius} \textquotedblleft{}a companion, an ally\textquotedblright{}.
\end{cousinweb}

\begin{grammarlens}[title={What you've built}]
One more word for telling, explaining and describing at B1.
\end{grammarlens}

\subsection*{Your turn}
\begin{itemize}
\raggedright
  \item \emph{Say it:} \emph{l'associazione}
  \item \emph{Say it:} \emph{l'associazione}, once more
  \item \emph{Say it:} a sentence of your own with \emph{l'associazione}
  \item \emph{Recall:} say the Italian for the voluntary work, then the Italian for the volunteer, then say \emph{l'associazione} again
\end{itemize}

\begin{tcolorbox}[breakable,colback=teal!4,colframe=teal!35!black,title={Before you move on}]
What does l'associazione mean? (\textquotedblleft{}The association, the organisation\textquotedblright{}.) Say it once more.
\end{tcolorbox}
\section[la beneficenza]{la beneficenza — the charity (giving)}

\label{lesson:IT-C367-beneficenza}



\begin{quote}
Before the new one: say the Italian for the volunteer, then the Italian for the association.
\end{quote}

\subsection*{You'll want to know: la beneficenza}
\textbf{la beneficenza} — \textquotedblleft{}the charity (giving)\textquotedblright{}.

\begin{quote}
I soldi del concerto vanno in beneficenza.
\end{quote}

\begin{quote}
\emph{The money from the concert goes to charity.}
\end{quote}

\emph{Dare in beneficenza} — to give to charity.

\begin{cousinweb}
\emph{beneficenza} is a learned borrowing from Latin \emph{beneficentia} \textquotedblleft{}kindness, doing good\textquotedblright{} (Lewis \& Short), built from \emph{bene} \textquotedblleft{}well\textquotedblright{} and \emph{facere} \textquotedblleft{}to do\textquotedblright{}.
\end{cousinweb}

\begin{grammarlens}[title={What you've built}]
One more word for telling, explaining and describing at B1.
\end{grammarlens}

\subsection*{Your turn}
\begin{itemize}
\raggedright
  \item \emph{Say it:} \emph{la beneficenza}
  \item \emph{Say it:} \emph{la beneficenza}, once more
  \item \emph{Say it:} a sentence of your own with \emph{la beneficenza}
  \item \emph{Recall:} say the Italian for the volunteer, then the Italian for the association, then say \emph{la beneficenza} again
\end{itemize}

\begin{tcolorbox}[breakable,colback=teal!4,colframe=teal!35!black,title={Before you move on}]
What does la beneficenza mean? (\textquotedblleft{}The charity (giving)\textquotedblright{}.) Say it once more.
\end{tcolorbox}
\section[la comunità]{la comunità — the community}

\label{lesson:IT-C367-comunita}



\begin{quote}
Before the new one: say the Italian for the association, then the Italian for the charity.
\end{quote}

\subsection*{You'll want to know: la comunità}
\textbf{la comunità} — \textquotedblleft{}the community\textquotedblright{}.

\begin{quote}
Tutta la comunità ha aiutato la famiglia.
\end{quote}

\begin{quote}
\emph{The whole community helped the family.}
\end{quote}

Also a residential centre: \emph{una comunità per ragazzi}.

\begin{cousinweb}
\emph{comunità} is a learned borrowing from Latin \emph{commūnitās} \textquotedblleft{}community, fellowship\textquotedblright{} (EtymDB), from \emph{commūnis} \textquotedblleft{}common, shared\textquotedblright{}.
\end{cousinweb}

\begin{grammarlens}[title={What you've built}]
One more word for telling, explaining and describing at B1.
\end{grammarlens}

\subsection*{Your turn}
\begin{itemize}
\raggedright
  \item \emph{Say it:} \emph{la comunità}
  \item \emph{Say it:} \emph{la comunità}, once more
  \item \emph{Say it:} a sentence of your own with \emph{la comunità}
  \item \emph{Recall:} say the Italian for the association, then the Italian for the charity, then say \emph{la comunità} again
\end{itemize}

\begin{tcolorbox}[breakable,colback=teal!4,colframe=teal!35!black,title={Before you move on}]
What does la comunità mean? (\textquotedblleft{}The community\textquotedblright{}.) Say it once more.
\end{tcolorbox}
